\documentclass[11pt,a4paper]{article}

% --- Codificación e idioma ---
\usepackage[utf8]{inputenc}
\usepackage[T1]{fontenc}
\usepackage[spanish,es-tabla]{babel}
\usepackage{lmodern}

% --- Geometría y formato general ---
\usepackage[margin=2.5cm]{geometry}
\usepackage{microtype}
\usepackage{parskip}
\usepackage{enumitem}
\setlist{topsep=2pt,itemsep=2pt,parsep=0pt}

% --- Tablas ---
\usepackage{array}
\usepackage{booktabs}

% --- Cabeceras y numeración ---
\usepackage{fancyhdr}
\pagestyle{fancy}
\fancyhf{}
\fancyhead[L]{\small Bases de Datos II}
\fancyhead[R]{\small Trabajo Práctico N.\textsuperscript{o} 3}
\fancyfoot[C]{\thepage}
\renewcommand{\headrulewidth}{0.4pt}

% --- Color y código ---
\usepackage{xcolor}
\definecolor{codebg}{RGB}{248,248,248}
\definecolor{codeframe}{RGB}{200,200,200}
\definecolor{codekw}{RGB}{0,0,180}
\definecolor{codestr}{RGB}{163,21,21}
\definecolor{codecmt}{RGB}{0,128,0}
\definecolor{notebg}{RGB}{255,250,230}
\definecolor{noteframe}{RGB}{220,180,80}

\usepackage{listings}
\lstdefinelanguage{SQLext}[]{SQL}{
    morekeywords={NUMERIC,VARCHAR,DECIMAL,DATE,TIMESTAMP,TIME,
        BIGINT,VARCHAR2,NUMBER,ENUM,SCHEMA,SEQUENCE,TRIGGER,
        BEFORE,AFTER,EACH,ROW,NEXTVAL,DUAL,RAISE,EXCEPTION,LOOP,
        IDENTIFIED,QUOTA,UNLIMITED,USERS,TABLESPACE,
        DATABASE,DOMAIN,GRANT,REVOKE,ROLE,USER,LOGIN,NOLOGIN,
        OPTION,REFERENCES,RESTRICT,CASCADE,
        AUTO\_INCREMENT,IF,EXISTS,ALTER,SHOW,USE,
        SEARCH\_PATH,EXECUTE,IMMEDIATE,PROCEDURE,FUNCTION,
        DELIMITER,SIGNAL,SQLSTATE,MESSAGE\_TEXT,
        DECLARE,BEGIN,END,RETURNS,LANGUAGE,PLPGSQL,
        OR,REPLACE,RETURN,FOR,IN,OUT,INOUT,
        WITH,FROM,TO,ON,AS,AND,OR,IS,NOT,NULL,
        INSERT,UPDATE,DELETE,SELECT,REVOKE,GRANT,
        SECURITY,INVOKER,DEFINER,ROUTINE,PROFILE,
        BIGSERIAL,SERIAL,IDENTITY,GENERATED,ALWAYS,
        DBMS\_OUTPUT,PUT\_LINE,RAISE\_APPLICATION\_ERROR,SYSDATE,
        SYSTEM,SYSDBA,DEFAULT,TEMPORARY},
    sensitive=false,
    morecomment=[l]{--},
    morestring=[b]'
}
\lstset{
    language=SQLext,
    basicstyle=\ttfamily\footnotesize,
    keywordstyle=\color{codekw}\bfseries,
    stringstyle=\color{codestr},
    commentstyle=\color{codecmt}\itshape,
    backgroundcolor=\color{codebg},
    frame=single,
    rulecolor=\color{codeframe},
    breaklines=true,
    breakatwhitespace=true,
    columns=fullflexible,
    keepspaces=true,
    showstringspaces=false,
    captionpos=b,
    xleftmargin=0.4em,
    xrightmargin=0.4em,
    aboveskip=8pt,
    belowskip=8pt,
    tabsize=2
}

% --- Cajas de nota ---
\usepackage[most]{tcolorbox}
\newtcolorbox{nota}{colback=notebg,colframe=noteframe,arc=2pt,
    left=6pt,right=6pt,top=4pt,bottom=4pt,boxrule=0.6pt}

% --- Hipervínculos ---
\usepackage[hidelinks,bookmarks=true]{hyperref}

\title{Trabajo Práctico N.\textsuperscript{o} 3:
       Procedimientos Almacenados y Triggers\\[4pt]
       \large Bases de Datos II}
\author{Tomás Rodeghiero}
\date{2026}

\begin{document}

% =====================================================
% Carátula
% =====================================================
\begin{titlepage}
    \centering
    \vspace*{1.5cm}

    {\large \textbf{Universidad Nacional de Río Cuarto}}\\[6pt]
    {\normalsize Facultad de Ciencias Exactas, Físico-Químicas y Naturales}\\[4pt]
    {\normalsize Departamento de Computación}\\[4pt]
    {\normalsize Licenciatura en Ciencias de la Computación}\\[3cm]

    {\Large \textbf{Bases de Datos II}}\\[1.2cm]

    {\Large \textbf{Trabajo Práctico N.\textsuperscript{o} 3}}\\[6pt]
    {\large Procedimientos Almacenados y Triggers}\\[3cm]

    \begin{tabular}{rl}
        \textbf{Alumno:} & Tomás Rodeghiero \\[4pt]
        \textbf{Año:}    & 2026 \\
    \end{tabular}

    \vfill
    {\small Resolución completa sobre los motores
    \emph{MySQL 8+}, \emph{PostgreSQL 14+} y \emph{Oracle XE},
    apoyada en las bases creadas en los Prácticos 1 y 2.}
\end{titlepage}

\tableofcontents
\newpage

% =====================================================
\section{Introducción}
% =====================================================

El Práctico 3 introduce dos mecanismos centrales del lado servidor:
\textbf{procedimientos y funciones almacenadas} y \textbf{triggers}.
Ambos permiten desplazar lógica de negocio desde la aplicación hasta el
motor de base de datos, garantizando reglas de integridad y de
auditoría que se cumplen \emph{siempre}, sin depender del cliente que
ejecute la operación.

Los ejercicios reutilizan las bases de datos del Práctico 1 (inciso 1.a
de facturación e inciso 1.b de estacionamiento) y, en el Ejercicio 4,
los usuarios definidos en el Práctico 2. Cada inciso indica qué motor
debe usarse:

\begin{itemize}
    \item Ejercicio 1: MySQL 8+ y PostgreSQL 14+.
    \item Ejercicio 2: PostgreSQL 14+ exclusivamente.
    \item Ejercicio 3: Oracle XE 10/11.
    \item Ejercicio 4: MySQL 8+ con el usuario \texttt{vendedor2}.
    \item Ejercicio 5: el motor en el que originalmente se creó cada
    base (MySQL/PostgreSQL para facturación, MySQL/Oracle para
    estacionamiento).
\end{itemize}

\begin{nota}
\textbf{Convenciones.} Los nombres de bases siguen los del Práctico 1:
\texttt{practico1a\_mysql}, \texttt{practico1a} (esquema en
PostgreSQL) y \texttt{practico1b\_mysql}. En MySQL los procedimientos
se rodean con \texttt{DELIMITER \$\$\dots\$\$ DELIMITER ;} para que el
cliente no corte la sentencia en cada \texttt{;} interno.
\end{nota}

% =====================================================
\section{Ejercicio 1: Cantidad máxima de ítems por factura}
% =====================================================

\paragraph{Enunciado.} Sobre la base de facturación del Práctico 1.a,
agregar a \texttt{factura} el campo \texttt{cantidad\_maxima\_items} y
garantizar que ninguna factura supere ese máximo.

\subsection*{Estrategia}

La consigna combina un cambio de estructura con una regla de
negocio que no se puede expresar como un \texttt{CHECK} declarativo
porque depende de \emph{otra tabla}. Por eso se resuelve con tres
elementos:

\begin{enumerate}
    \item Un \texttt{ALTER TABLE} que agrega la columna y un
    \texttt{CHECK} sobre su rango.
    \item Un trigger \texttt{BEFORE INSERT} (y \texttt{BEFORE UPDATE}
    cuando un ítem se reasigna a otra factura) sobre
    \texttt{item\_factura}, que cuenta los ítems existentes y aborta
    si la inserción supera el límite.
    \item Un trigger \texttt{BEFORE UPDATE} sobre \texttt{factura} que
    impide bajar \texttt{cantidad\_maxima\_items} por debajo de los
    ítems ya cargados.
\end{enumerate}

\subsection{MySQL 8+}

\begin{lstlisting}
USE practico1a_mysql;

-- 1) Agregar la columna con un default razonable y validacion de rango.
ALTER TABLE factura
    ADD COLUMN cantidad_maxima_items INT NOT NULL DEFAULT 10,
    ADD CONSTRAINT ck_factura_cantidad_maxima_items
        CHECK (cantidad_maxima_items > 0);

-- 2) Trigger BEFORE INSERT en item_factura.
DROP TRIGGER IF EXISTS bi_item_factura_control_max_items;
DELIMITER $$
CREATE TRIGGER bi_item_factura_control_max_items
BEFORE INSERT ON item_factura
FOR EACH ROW
BEGIN
    DECLARE v_max INT;
    DECLARE v_actual INT;

    SELECT f.cantidad_maxima_items INTO v_max
      FROM factura f
     WHERE f.nro_factura = NEW.nro_factura;

    SELECT COUNT(*) INTO v_actual
      FROM item_factura i
     WHERE i.nro_factura = NEW.nro_factura;

    IF v_actual + 1 > v_max THEN
        SIGNAL SQLSTATE '45000'
            SET MESSAGE_TEXT = 'La factura supera la cantidad maxima de items permitidos';
    END IF;
END$$
DELIMITER ;

-- 3) Trigger BEFORE UPDATE en item_factura: solo valida el cambio de
--    factura destino (mover un item de una factura a otra).
DROP TRIGGER IF EXISTS bu_item_factura_control_max_items;
DELIMITER $$
CREATE TRIGGER bu_item_factura_control_max_items
BEFORE UPDATE ON item_factura
FOR EACH ROW
BEGIN
    DECLARE v_max INT;
    DECLARE v_actual INT;

    IF NEW.nro_factura <> OLD.nro_factura THEN
        SELECT f.cantidad_maxima_items INTO v_max
          FROM factura f
         WHERE f.nro_factura = NEW.nro_factura;

        SELECT COUNT(*) INTO v_actual
          FROM item_factura i
         WHERE i.nro_factura = NEW.nro_factura;

        IF v_actual + 1 > v_max THEN
            SIGNAL SQLSTATE '45000'
                SET MESSAGE_TEXT = 'La factura destino supera la cantidad maxima de items permitidos';
        END IF;
    END IF;
END$$
DELIMITER ;

-- 4) Trigger BEFORE UPDATE en factura: impide bajar el maximo por
--    debajo de los items ya cargados.
DROP TRIGGER IF EXISTS bu_factura_control_max_items;
DELIMITER $$
CREATE TRIGGER bu_factura_control_max_items
BEFORE UPDATE ON factura
FOR EACH ROW
BEGIN
    DECLARE v_actual INT;

    IF NEW.cantidad_maxima_items <= 0 THEN
        SIGNAL SQLSTATE '45000'
            SET MESSAGE_TEXT = 'cantidad_maxima_items debe ser mayor que 0';
    END IF;

    IF NEW.cantidad_maxima_items <> OLD.cantidad_maxima_items THEN
        SELECT COUNT(*) INTO v_actual
          FROM item_factura i
         WHERE i.nro_factura = OLD.nro_factura;

        IF v_actual > NEW.cantidad_maxima_items THEN
            SIGNAL SQLSTATE '45000'
                SET MESSAGE_TEXT = 'No se puede reducir el maximo: ya hay mas items cargados';
        END IF;
    END IF;
END$$
DELIMITER ;
\end{lstlisting}

\paragraph{¿Por qué tres triggers?} El trigger de
\texttt{INSERT} cubre la altas; el de \texttt{UPDATE} sobre
\texttt{item\_factura} sólo se activa cuando el ítem cambia de
factura, evitando trabajo innecesario en updates de
\texttt{cantidad}/\texttt{precio}; y el trigger de \texttt{UPDATE}
sobre \texttt{factura} cubre el caso simétrico, en el que el problema
no es el ítem que entra sino el límite que baja.

\subsection{PostgreSQL 14+}

PostgreSQL separa la \emph{función} (que contiene la lógica) del
\emph{trigger} (que la dispara). Una sola función cubre
\texttt{INSERT} y \texttt{UPDATE} usando \texttt{TG\_OP}.

\begin{lstlisting}
SET search_path TO practico1a;

-- 1) Agregar columna con default y CHECK.
ALTER TABLE practico1a.factura
    ADD COLUMN cantidad_maxima_items INTEGER NOT NULL DEFAULT 10;

ALTER TABLE practico1a.factura
    ADD CONSTRAINT ck_factura_cantidad_maxima_items
    CHECK (cantidad_maxima_items > 0);

-- 2) Funcion + trigger de control sobre item_factura.
DROP TRIGGER IF EXISTS biu_item_factura_control_max_items
    ON practico1a.item_factura;
DROP FUNCTION IF EXISTS practico1a.fn_control_max_items_item_factura();

CREATE OR REPLACE FUNCTION practico1a.fn_control_max_items_item_factura()
RETURNS TRIGGER
LANGUAGE plpgsql
AS $$
DECLARE
    v_max INTEGER;
    v_actual INTEGER;
BEGIN
    -- En UPDATE de la misma factura no cambia la cantidad de items.
    IF TG_OP = 'UPDATE' AND NEW.nro_factura = OLD.nro_factura THEN
        RETURN NEW;
    END IF;

    SELECT f.cantidad_maxima_items INTO v_max
      FROM practico1a.factura f
     WHERE f.nro_factura = NEW.nro_factura
     FOR UPDATE;

    SELECT COUNT(*) INTO v_actual
      FROM practico1a.item_factura i
     WHERE i.nro_factura = NEW.nro_factura;

    IF v_actual + 1 > v_max THEN
        RAISE EXCEPTION
          'La factura % supera la cantidad maxima de items permitidos (%).',
          NEW.nro_factura, v_max;
    END IF;

    RETURN NEW;
END;
$$;

CREATE TRIGGER biu_item_factura_control_max_items
BEFORE INSERT OR UPDATE ON practico1a.item_factura
FOR EACH ROW
EXECUTE FUNCTION practico1a.fn_control_max_items_item_factura();

-- 3) Funcion + trigger de control sobre factura.cantidad_maxima_items.
DROP TRIGGER IF EXISTS bu_factura_control_max_items
    ON practico1a.factura;
DROP FUNCTION IF EXISTS practico1a.fn_control_update_max_items_factura();

CREATE OR REPLACE FUNCTION practico1a.fn_control_update_max_items_factura()
RETURNS TRIGGER
LANGUAGE plpgsql
AS $$
DECLARE
    v_actual INTEGER;
BEGIN
    IF NEW.cantidad_maxima_items <= 0 THEN
        RAISE EXCEPTION 'cantidad_maxima_items debe ser mayor que 0';
    END IF;

    IF NEW.cantidad_maxima_items <> OLD.cantidad_maxima_items THEN
        SELECT COUNT(*) INTO v_actual
          FROM practico1a.item_factura i
         WHERE i.nro_factura = OLD.nro_factura;

        IF v_actual > NEW.cantidad_maxima_items THEN
            RAISE EXCEPTION
              'No se puede reducir el maximo de la factura %: tiene % items cargados',
              OLD.nro_factura, v_actual;
        END IF;
    END IF;

    RETURN NEW;
END;
$$;

CREATE TRIGGER bu_factura_control_max_items
BEFORE UPDATE OF cantidad_maxima_items ON practico1a.factura
FOR EACH ROW
EXECUTE FUNCTION practico1a.fn_control_update_max_items_factura();
\end{lstlisting}

\paragraph{Detalles a destacar.}

\begin{itemize}
    \item En PostgreSQL se usa \texttt{FOR UPDATE} sobre la fila de
    \texttt{factura} para evitar condiciones de carrera bajo
    concurrencia: si dos sesiones intentan insertar el último ítem al
    mismo tiempo, una espera a que la otra confirme.
    \item El trigger \texttt{BEFORE UPDATE OF cantidad\_maxima\_items}
    sólo dispara cuando esa columna cambia, lo que es más eficiente
    que disparar en cualquier \texttt{UPDATE}.
\end{itemize}

\subsection*{Prueba (válida para ambos motores)}

\begin{lstlisting}
-- Suponiendo que la factura 1001 ya tiene 2 items cargados.
UPDATE factura SET cantidad_maxima_items = 2 WHERE nro_factura = 1001;

-- FALLA: dejaria 3 items en la factura 1001.
INSERT INTO item_factura (cod_producto, nro_factura, cantidad, precio)
VALUES (101, 1001, 1, 15000.00);

-- FALLA: intenta reducir el maximo por debajo de los 2 items existentes.
UPDATE factura SET cantidad_maxima_items = 1 WHERE nro_factura = 1001;

-- FUNCIONA: subir el maximo siempre es valido.
UPDATE factura SET cantidad_maxima_items = 5 WHERE nro_factura = 1001;
\end{lstlisting}

% =====================================================
\section{Ejercicio 2: Cotización máxima por banco (PostgreSQL)}
% =====================================================

\paragraph{Enunciado.} Crear las tablas \texttt{Banco} y
\texttt{Cotizacion}, e implementar una función que reciba un banco,
calcule la cotización máxima \emph{sin usar} \texttt{MAX}, actualice
\texttt{Banco.cotizacion\_maxima} y devuelva el resultado por un
parámetro \texttt{OUT}.

\subsection*{Diseño}

Como la consigna prohíbe \texttt{MAX}, la función recorre las
cotizaciones del banco con un \texttt{FOR \dots\ LOOP} en PL/pgSQL y
mantiene una variable acumuladora con el mayor valor visto hasta el
momento. Si el banco no tiene cotizaciones cargadas o no existe, la
función levanta una excepción explícita.

\begin{lstlisting}
CREATE SCHEMA IF NOT EXISTS practico3_ej2;
SET search_path TO practico3_ej2;

DROP TABLE IF EXISTS cotizacion;
DROP TABLE IF EXISTS banco;

CREATE TABLE banco (
    cod_banco INTEGER PRIMARY KEY,
    nombre VARCHAR(80) NOT NULL,
    precio NUMERIC(12,4) NOT NULL,
    cotizacion_maxima NUMERIC(12,4),
    CONSTRAINT ck_banco_precio CHECK (precio > 0),
    CONSTRAINT ck_banco_cot_max
        CHECK (cotizacion_maxima IS NULL OR cotizacion_maxima > 0)
);

CREATE TABLE cotizacion (
    cod_cotizacion INTEGER GENERATED ALWAYS AS IDENTITY PRIMARY KEY,
    cod_banco INTEGER NOT NULL,
    cotizacion NUMERIC(12,4) NOT NULL,
    CONSTRAINT ck_cotizacion_valor CHECK (cotizacion > 0),
    CONSTRAINT fk_cotizacion_banco
        FOREIGN KEY (cod_banco) REFERENCES banco(cod_banco)
        ON DELETE RESTRICT
);

CREATE INDEX idx_cotizacion_cod_banco ON cotizacion(cod_banco);
\end{lstlisting}

\subsection*{Función con parámetros \texttt{IN}/\texttt{OUT}}

\begin{lstlisting}
CREATE OR REPLACE FUNCTION practico3_ej2.fn_actualizar_cotizacion_maxima(
    IN  p_cod_banco INTEGER,
    OUT p_cotizacion_maxima NUMERIC(12,4)
)
LANGUAGE plpgsql
AS $$
DECLARE
    v_cot NUMERIC(12,4);
BEGIN
    -- Verificacion de existencia + bloqueo del banco.
    PERFORM 1
      FROM practico3_ej2.banco b
     WHERE b.cod_banco = p_cod_banco
     FOR UPDATE;

    IF NOT FOUND THEN
        RAISE EXCEPTION 'No existe banco con cod_banco=%', p_cod_banco;
    END IF;

    p_cotizacion_maxima := NULL;

    -- Calculo manual del maximo, sin MAX.
    FOR v_cot IN
        SELECT c.cotizacion
          FROM practico3_ej2.cotizacion c
         WHERE c.cod_banco = p_cod_banco
    LOOP
        IF p_cotizacion_maxima IS NULL OR v_cot > p_cotizacion_maxima THEN
            p_cotizacion_maxima := v_cot;
        END IF;
    END LOOP;

    IF p_cotizacion_maxima IS NULL THEN
        RAISE EXCEPTION 'El banco % no tiene cotizaciones cargadas', p_cod_banco;
    END IF;

    -- Persistir el resultado en la tabla banco.
    UPDATE practico3_ej2.banco b
       SET cotizacion_maxima = p_cotizacion_maxima
     WHERE b.cod_banco = p_cod_banco;

    RETURN;
END;
$$;
\end{lstlisting}

\paragraph{Análisis.}

\begin{itemize}
    \item \textbf{Parámetros}: \texttt{p\_cod\_banco} entra como
    \texttt{IN}, \texttt{p\_cotizacion\_maxima} sale como \texttt{OUT}.
    PostgreSQL trata las funciones con \texttt{OUT} como devolviendo
    una fila con esas columnas; por eso se invoca con
    \texttt{SELECT * FROM fn\dots(\dots)}.
    \item \textbf{Estructura iterativa}: el \texttt{FOR \dots\ LOOP}
    avanza una fila por iteración. La variable
    \texttt{p\_cotizacion\_maxima} se inicializa en \texttt{NULL} y
    se actualiza con cada valor mayor.
    \item \textbf{Manejo de errores}: dos \texttt{RAISE EXCEPTION}
    distintos para los dos casos de borde (banco inexistente y banco
    sin cotizaciones), con mensajes diferentes para que el cliente
    pueda diferenciarlos.
    \item \textbf{Concurrencia}: \texttt{PERFORM \dots\ FOR UPDATE}
    bloquea la fila del banco para que dos invocaciones simultáneas
    no se pisen al actualizar \texttt{cotizacion\_maxima}.
\end{itemize}

\subsection*{Datos de prueba y ejecución}

\begin{lstlisting}
SET search_path TO practico3_ej2;

INSERT INTO banco (cod_banco, nombre, precio, cotizacion_maxima) VALUES
    (1, 'Banco Rio', 1200.0000, NULL),
    (2, 'Banco Sur',  950.0000, NULL);

INSERT INTO cotizacion (cod_banco, cotizacion) VALUES
    (1, 1210.5000), (1, 1198.2000), (1, 1245.7500), (1, 1222.1000),
    (2,  960.0000), (2, 1001.3000), (2,  998.9000);

-- Llamado: la salida es 1245.7500 para el banco 1, 1001.3000 para el 2.
SELECT * FROM practico3_ej2.fn_actualizar_cotizacion_maxima(1);
SELECT * FROM practico3_ej2.fn_actualizar_cotizacion_maxima(2);

SELECT cod_banco, nombre, cotizacion_maxima
  FROM practico3_ej2.banco
 ORDER BY cod_banco;
\end{lstlisting}

% =====================================================
\section{Ejercicio 3: Saldo de cuenta bancaria (Oracle)}
% =====================================================

\paragraph{Enunciado.} Crear las tablas \texttt{cuenta} y
\texttt{movimiento} y dos procedimientos: (a) registrar un movimiento
manteniendo el saldo y mostrarlo por pantalla, y (b) calcular el saldo
de una cuenta a una fecha dada, mostrándolo por pantalla y
devolviéndolo por un parámetro de salida.

\subsection*{Esquema}

El saldo de cada cuenta se mantiene \emph{denormalizado} en
\texttt{cuenta.saldo} para evitar recalcularlo a partir de toda la
historia en cada consulta. La consistencia entre ambos se mantiene
mediante el procedimiento de inserción.

\begin{lstlisting}
SET SERVEROUTPUT ON;

-- Tablas
CREATE TABLE cuenta (
    nro_cuenta NUMBER(10) CONSTRAINT pk_cuenta PRIMARY KEY,
    saldo NUMBER(14,2) DEFAULT 0 NOT NULL
);

CREATE TABLE movimiento (
    nro_movimiento NUMBER(12) CONSTRAINT pk_movimiento PRIMARY KEY,
    nro_cuenta NUMBER(10) NOT NULL,
    fecha DATE DEFAULT SYSDATE NOT NULL,
    debe NUMBER(14,2) DEFAULT 0 NOT NULL,
    haber NUMBER(14,2) DEFAULT 0 NOT NULL,
    CONSTRAINT fk_movimiento_cuenta
        FOREIGN KEY (nro_cuenta) REFERENCES cuenta(nro_cuenta),
    CONSTRAINT ck_mov_montos_no_negativos
        CHECK (debe >= 0 AND haber >= 0),
    CONSTRAINT ck_mov_operacion_valida
        CHECK ((debe > 0 AND haber = 0) OR (haber > 0 AND debe = 0))
);

-- En Oracle 10g/11g no hay IDENTITY: secuencia + trigger.
CREATE SEQUENCE seq_movimiento
    START WITH 1 INCREMENT BY 1 NOCACHE NOCYCLE;

CREATE OR REPLACE TRIGGER bi_movimiento
BEFORE INSERT ON movimiento
FOR EACH ROW
BEGIN
    IF :NEW.nro_movimiento IS NULL THEN
        SELECT seq_movimiento.NEXTVAL INTO :NEW.nro_movimiento FROM dual;
    END IF;

    IF :NEW.fecha IS NULL THEN
        :NEW.fecha := SYSDATE;
    END IF;
END;
/
\end{lstlisting}

\subsection{Inciso a) Procedimiento que inserta movimiento y mantiene saldo}

\begin{lstlisting}
CREATE OR REPLACE PROCEDURE sp_registrar_movimiento (
    p_nro_cuenta IN cuenta.nro_cuenta%TYPE,
    p_debe       IN NUMBER,
    p_haber      IN NUMBER
)
IS
    v_saldo_actual cuenta.saldo%TYPE;
    v_debe  NUMBER(14,2) := NVL(p_debe, 0);
    v_haber NUMBER(14,2) := NVL(p_haber, 0);
BEGIN
    -- Validaciones de los parametros.
    IF v_debe < 0 OR v_haber < 0 THEN
        RAISE_APPLICATION_ERROR(-20001,
            'Debe/Haber no pueden ser negativos');
    END IF;

    IF NOT ((v_debe > 0 AND v_haber = 0)
         OR (v_haber > 0 AND v_debe = 0)) THEN
        RAISE_APPLICATION_ERROR(-20002,
            'La operacion debe tener solo DEBE o solo HABER');
    END IF;

    -- Bloqueo de la cuenta para mantener consistencia bajo concurrencia.
    SELECT c.saldo INTO v_saldo_actual
      FROM cuenta c
     WHERE c.nro_cuenta = p_nro_cuenta
     FOR UPDATE;

    INSERT INTO movimiento (nro_cuenta, fecha, debe, haber)
    VALUES (p_nro_cuenta, SYSDATE, v_debe, v_haber);

    UPDATE cuenta c
       SET c.saldo = c.saldo + v_debe - v_haber
     WHERE c.nro_cuenta = p_nro_cuenta
     RETURNING c.saldo INTO v_saldo_actual;

    DBMS_OUTPUT.PUT_LINE(
        'Cuenta: ' || p_nro_cuenta || ' | Saldo actual: ' ||
        TO_CHAR(v_saldo_actual, 'FM9999999990D00')
    );
EXCEPTION
    WHEN NO_DATA_FOUND THEN
        RAISE_APPLICATION_ERROR(-20003,
            'No existe la cuenta ' || p_nro_cuenta);
END;
/
\end{lstlisting}

\paragraph{Lógica del procedimiento.}

\begin{enumerate}
    \item Recibe los parámetros \texttt{IN}: \texttt{p\_nro\_cuenta},
    \texttt{p\_debe}, \texttt{p\_haber}.
    \item Reemplaza eventuales \texttt{NULL} por \texttt{0} con
    \texttt{NVL} y valida que la operación sea \emph{sólo} debe o
    \emph{sólo} haber.
    \item Bloquea la fila de la cuenta con \texttt{SELECT \dots\ FOR
    UPDATE} para evitar condiciones de carrera entre dos
    movimientos concurrentes sobre la misma cuenta.
    \item Inserta el movimiento (la columna \texttt{nro\_movimiento}
    se completa automáticamente con el trigger
    \texttt{bi\_movimiento}).
    \item Actualiza el saldo en la misma transacción y captura el
    nuevo valor con \texttt{RETURNING}.
    \item Imprime el resultado con \texttt{DBMS\_OUTPUT.PUT\_LINE}.
    \item En caso de cuenta inexistente, atrapa
    \texttt{NO\_DATA\_FOUND} y emite un error de aplicación
    diferenciado.
\end{enumerate}

\subsection{Inciso b) Saldo a una fecha dada}

\begin{lstlisting}
CREATE OR REPLACE PROCEDURE sp_saldo_a_fecha (
    p_nro_cuenta IN  cuenta.nro_cuenta%TYPE,
    p_fecha      IN  DATE,
    p_saldo      OUT NUMBER
)
IS
    v_saldo_actual    cuenta.saldo%TYPE;
    v_delta_posterior NUMBER(14,2);
BEGIN
    SELECT c.saldo INTO v_saldo_actual
      FROM cuenta c
     WHERE c.nro_cuenta = p_nro_cuenta;

    -- Suma del impacto de los movimientos POSTERIORES a la fecha.
    SELECT NVL(SUM(m.debe - m.haber), 0)
      INTO v_delta_posterior
      FROM movimiento m
     WHERE m.nro_cuenta = p_nro_cuenta
       AND m.fecha > p_fecha;

    p_saldo := v_saldo_actual - v_delta_posterior;

    DBMS_OUTPUT.PUT_LINE(
        'Cuenta: ' || p_nro_cuenta ||
        ' | Fecha: ' || TO_CHAR(p_fecha, 'YYYY-MM-DD HH24:MI:SS') ||
        ' | Saldo: ' || TO_CHAR(p_saldo, 'FM9999999990D00')
    );
EXCEPTION
    WHEN NO_DATA_FOUND THEN
        RAISE_APPLICATION_ERROR(-20004,
            'No existe la cuenta ' || p_nro_cuenta);
END;
/
\end{lstlisting}

\paragraph{Idea del cálculo.} Se parte del saldo actual y se le resta
el efecto neto de los movimientos cuya fecha sea estrictamente mayor
a \texttt{p\_fecha}: así se obtiene cuánto valía la cuenta al
\emph{cierre} de \texttt{p\_fecha}. La alternativa --- sumar todos los
movimientos hasta esa fecha desde el comienzo --- es equivalente pero
más cara cuando hay muchos movimientos antiguos.

\subsection*{Prueba}

\begin{lstlisting}
SET SERVEROUTPUT ON;

INSERT INTO cuenta (nro_cuenta, saldo) VALUES (1001, 0);
INSERT INTO cuenta (nro_cuenta, saldo) VALUES (1002, 500);

BEGIN
    sp_registrar_movimiento(1001, 1000, 0);  -- saldo 1000
    sp_registrar_movimiento(1001, 0, 150);   -- saldo  850
    sp_registrar_movimiento(1001, 200, 0);   -- saldo 1050
END;
/

DECLARE
    v_saldo NUMBER;
BEGIN
    sp_saldo_a_fecha(1001, SYSDATE, v_saldo);
    DBMS_OUTPUT.PUT_LINE(
        'Saldo retornado por OUT: ' ||
        TO_CHAR(v_saldo, 'FM9999999990D00'));
END;
/
\end{lstlisting}

% =====================================================
\section{Ejercicio 4: Detección de facturas inconsistentes (MySQL)}
% =====================================================

\paragraph{Enunciado.} Crear un procedimiento, en MySQL, que
detecte facturas \emph{inconsistentes} (cuando la suma de
\texttt{cantidad $\times$ precio} de sus ítems no coincide con
\texttt{factura.monto}) y las almacene en una tabla. El procedimiento
debe ser creado con el usuario \texttt{vendedor2} del Práctico 2 y
ejecutarse con permisos del \emph{invocador}.

\subsection*{Preparación (como \texttt{root})}

En el Práctico 2 (inciso 1.f) se le revocaron todos los privilegios a
\texttt{vendedor2}. Antes de poder crear el procedimiento como
\texttt{vendedor2}, hay que reotorgarle los permisos mínimos y, dado
que el procedimiento usará \texttt{SQL SECURITY INVOKER}, también hay
que dotar de los privilegios necesarios a \texttt{vendedor1} (el otro
usuario que lo invocará).

\begin{lstlisting}
USE practico1a_mysql;

-- Tabla destino donde se guardan las inconsistencias detectadas.
CREATE TABLE IF NOT EXISTS facturas_inconsistentes (
    nro_factura       INT PRIMARY KEY,
    monto_factura     DECIMAL(12,2) NOT NULL,
    monto_items       DECIMAL(12,2) NOT NULL,
    diferencia        DECIMAL(12,2) NOT NULL,
    fecha_control     DATETIME NOT NULL,
    usuario_ejecucion VARCHAR(100) NOT NULL
);

-- Permiso minimo para que vendedor2 pueda crear procedimientos.
GRANT CREATE ROUTINE ON practico1a_mysql.*
    TO 'vendedor2'@'localhost';

-- Como el procedimiento corre con SQL SECURITY INVOKER,
-- cada usuario que lo ejecute necesita estos permisos:
GRANT SELECT (nro_factura, monto)
   ON practico1a_mysql.factura
   TO 'vendedor1'@'localhost', 'vendedor2'@'localhost';

GRANT SELECT (nro_factura, cantidad, precio)
   ON practico1a_mysql.item_factura
   TO 'vendedor1'@'localhost', 'vendedor2'@'localhost';

GRANT INSERT, UPDATE, SELECT
   ON practico1a_mysql.facturas_inconsistentes
   TO 'vendedor1'@'localhost', 'vendedor2'@'localhost';
\end{lstlisting}

\subsection*{Creación del procedimiento (conectado como
\texttt{vendedor2})}

\begin{lstlisting}
-- $ mysql -u vendedor2 -p practico1a_mysql

DROP PROCEDURE IF EXISTS sp_guardar_facturas_inconsistentes;
DELIMITER $$
CREATE PROCEDURE sp_guardar_facturas_inconsistentes()
SQL SECURITY INVOKER
BEGIN
    INSERT INTO facturas_inconsistentes (
        nro_factura, monto_factura, monto_items,
        diferencia, fecha_control, usuario_ejecucion
    )
    SELECT
        f.nro_factura,
        f.monto AS monto_factura,
        ROUND(COALESCE(SUM(i.cantidad * i.precio), 0), 2) AS monto_items,
        ROUND(COALESCE(SUM(i.cantidad * i.precio), 0) - f.monto, 2)
            AS diferencia,
        NOW()  AS fecha_control,
        USER() AS usuario_ejecucion
      FROM factura f
      LEFT JOIN item_factura i ON i.nro_factura = f.nro_factura
     GROUP BY f.nro_factura, f.monto
    HAVING ROUND(COALESCE(SUM(i.cantidad * i.precio), 0), 2)
        <> ROUND(f.monto, 2)
    ON DUPLICATE KEY UPDATE
        monto_factura     = VALUES(monto_factura),
        monto_items       = VALUES(monto_items),
        diferencia        = VALUES(diferencia),
        fecha_control     = VALUES(fecha_control),
        usuario_ejecucion = VALUES(usuario_ejecucion);
END$$
DELIMITER ;
\end{lstlisting}

\paragraph{Por qué \texttt{SQL SECURITY INVOKER}.} Por defecto MySQL
crea procedimientos con \texttt{SQL SECURITY DEFINER}, lo que hace que
todas las ejecuciones se realicen con los permisos de quien lo creó.
La consigna pide lo contrario: que el procedimiento corra con los
permisos de quien lo invoca. Eso obliga a que \emph{cada} invocador
posea los permisos que la rutina necesita; de ahí los \texttt{GRANT}
de la sección anterior.

\paragraph{Verificación del \texttt{DEFINER}.} Para confirmar que se
creó como \texttt{vendedor2} y con \texttt{SQL SECURITY INVOKER}:

\begin{lstlisting}
SHOW CREATE PROCEDURE sp_guardar_facturas_inconsistentes\G
-- En la salida debe verse:
--   DEFINER = `vendedor2`@`localhost`
--   SQL SECURITY INVOKER
\end{lstlisting}

\subsection*{Permiso de ejecución (como \texttt{root})}

\begin{lstlisting}
GRANT EXECUTE ON PROCEDURE practico1a_mysql.sp_guardar_facturas_inconsistentes
    TO 'vendedor1'@'localhost', 'vendedor2'@'localhost';
\end{lstlisting}

\subsection*{Prueba con ambos usuarios}

Para forzar una inconsistencia (sólo \texttt{root} suele tener
\texttt{UPDATE} sobre \texttt{factura}):

\begin{lstlisting}
USE practico1a_mysql;
UPDATE factura SET monto = monto + 100 WHERE nro_factura = 1001;
\end{lstlisting}

Luego, conectándose como cada usuario y ejecutando:

\begin{lstlisting}
-- $ mysql -u vendedor2 -p practico1a_mysql
CALL sp_guardar_facturas_inconsistentes();
SELECT * FROM facturas_inconsistentes
 ORDER BY fecha_control DESC, nro_factura;

-- $ mysql -u vendedor1 -p practico1a_mysql
CALL sp_guardar_facturas_inconsistentes();
SELECT * FROM facturas_inconsistentes
 ORDER BY fecha_control DESC, nro_factura;
\end{lstlisting}

\paragraph{¿Por qué \texttt{ON DUPLICATE KEY UPDATE}?} Porque
\texttt{nro\_factura} es \texttt{PRIMARY KEY} en
\texttt{facturas\_inconsistentes}: si una factura ya estaba marcada
como inconsistente y se vuelve a ejecutar el procedimiento, en lugar
de fallar por clave duplicada se actualiza con la última lectura. La
columna \texttt{usuario\_ejecucion} usa \texttt{USER()}, que en MySQL
devuelve el usuario realmente conectado, lo que permite distinguir
quién corrió el control.

% =====================================================
\section{Ejercicio 5: Triggers y autonumerado}
% =====================================================

\subsection{a) Mayúsculas antes de insertar un producto}

Se trata de un trigger \texttt{BEFORE INSERT} que normaliza el dato
antes de guardarlo, de modo que la auditoría posterior y las búsquedas
queden uniformes.

\subsubsection*{MySQL 8+ (\texttt{practico1a\_mysql})}

\begin{lstlisting}
USE practico1a_mysql;

DROP TRIGGER IF EXISTS bi_producto_mayusculas;
DELIMITER $$
CREATE TRIGGER bi_producto_mayusculas
BEFORE INSERT ON producto
FOR EACH ROW
BEGIN
    SET NEW.descripcion = UPPER(TRIM(NEW.descripcion));
END$$
DELIMITER ;
\end{lstlisting}

\subsubsection*{PostgreSQL 14+ (\texttt{practico1a})}

\begin{lstlisting}
SET search_path TO practico1a;

DROP TRIGGER IF EXISTS bi_producto_mayusculas ON practico1a.producto;
DROP FUNCTION IF EXISTS practico1a.fn_producto_mayusculas();

CREATE OR REPLACE FUNCTION practico1a.fn_producto_mayusculas()
RETURNS TRIGGER
LANGUAGE plpgsql
AS $$
BEGIN
    NEW.descripcion := UPPER(BTRIM(NEW.descripcion));
    RETURN NEW;
END;
$$;

CREATE TRIGGER bi_producto_mayusculas
BEFORE INSERT ON practico1a.producto
FOR EACH ROW
EXECUTE FUNCTION practico1a.fn_producto_mayusculas();
\end{lstlisting}

\paragraph{Detalle.} En MySQL se asigna \texttt{NEW.descripcion} con
\texttt{SET}; en PostgreSQL se asigna directamente con \texttt{:=} y
se devuelve \texttt{NEW}. El \texttt{TRIM/BTRIM} adicional limpia
espacios al principio y al final, defendiéndose de errores de
formato.

\subsection{b) Auditoría de cambios en \texttt{producto.cantidad}}

La auditoría debe registrar \texttt{cod\_producto}, el
\emph{movimiento} (\texttt{NEW.cantidad - OLD.cantidad}), la fecha y
el \emph{usuario logueado en la aplicación cliente}. Como el motor no
tiene forma de conocer el usuario de aplicación por sí mismo, la
aplicación lo provee a través de una variable de sesión que el
trigger lee.

\subsubsection*{MySQL 8+}

\begin{lstlisting}
USE practico1a_mysql;

CREATE TABLE IF NOT EXISTS auditoriaProducto (
    id_auditoria        BIGINT AUTO_INCREMENT PRIMARY KEY,
    cod_producto        INT NOT NULL,
    movimiento          INT NOT NULL,
    fecha_actualizacion DATETIME NOT NULL,
    usuario_aplicacion  VARCHAR(100) NOT NULL,
    CONSTRAINT fk_auditoria_producto
        FOREIGN KEY (cod_producto) REFERENCES producto(cod_producto)
);

DROP TRIGGER IF EXISTS au_producto_auditoria_cantidad;
DELIMITER $$
CREATE TRIGGER au_producto_auditoria_cantidad
AFTER UPDATE ON producto
FOR EACH ROW
BEGIN
    DECLARE v_usuario_app VARCHAR(100);

    IF NEW.cantidad <> OLD.cantidad THEN
        -- La aplicacion debe setear esta variable por sesion.
        --   SET @app_user_code = 'VENDEDOR_APP_01';
        SET v_usuario_app = @app_user_code;

        IF v_usuario_app IS NULL OR TRIM(v_usuario_app) = '' THEN
            SIGNAL SQLSTATE '45000'
                SET MESSAGE_TEXT =
                  'Debe definir @app_user_code en la sesion para auditar';
        END IF;

        INSERT INTO auditoriaProducto (
            cod_producto, movimiento,
            fecha_actualizacion, usuario_aplicacion
        )
        VALUES (
            NEW.cod_producto,
            NEW.cantidad - OLD.cantidad,
            NOW(),
            v_usuario_app
        );
    END IF;
END$$
DELIMITER ;
\end{lstlisting}

\paragraph{Uso desde la aplicación.}

\begin{lstlisting}
SET @app_user_code = 'VENDEDOR_APP_01';
UPDATE producto SET cantidad = cantidad + 5 WHERE cod_producto = 101;

SELECT * FROM auditoriaProducto ORDER BY id_auditoria DESC;
\end{lstlisting}

\subsubsection*{PostgreSQL 14+}

\begin{lstlisting}
SET search_path TO practico1a;

CREATE TABLE IF NOT EXISTS practico1a.auditoria_producto (
    id_auditoria        BIGSERIAL PRIMARY KEY,
    cod_producto        INTEGER NOT NULL
        REFERENCES practico1a.producto(cod_producto),
    movimiento          INTEGER NOT NULL,
    fecha_actualizacion TIMESTAMP NOT NULL,
    usuario_aplicacion  VARCHAR(100) NOT NULL
);

DROP TRIGGER IF EXISTS au_producto_auditoria_cantidad
    ON practico1a.producto;
DROP FUNCTION IF EXISTS practico1a.fn_auditoria_producto_cantidad();

CREATE OR REPLACE FUNCTION practico1a.fn_auditoria_producto_cantidad()
RETURNS TRIGGER
LANGUAGE plpgsql
AS $$
DECLARE
    v_usuario_app TEXT;
BEGIN
    IF NEW.cantidad <> OLD.cantidad THEN
        -- La aplicacion debe setear esta variable de sesion.
        --   SET app.user_code = 'VENDEDOR_APP_01';
        v_usuario_app := current_setting('app.user_code', true);

        IF v_usuario_app IS NULL OR BTRIM(v_usuario_app) = '' THEN
            RAISE EXCEPTION
              'Debe definir app.user_code en la sesion para auditar';
        END IF;

        INSERT INTO practico1a.auditoria_producto (
            cod_producto, movimiento,
            fecha_actualizacion, usuario_aplicacion
        )
        VALUES (
            NEW.cod_producto,
            NEW.cantidad - OLD.cantidad,
            NOW(),
            v_usuario_app
        );
    END IF;

    RETURN NEW;
END;
$$;

CREATE TRIGGER au_producto_auditoria_cantidad
AFTER UPDATE OF cantidad ON practico1a.producto
FOR EACH ROW
EXECUTE FUNCTION practico1a.fn_auditoria_producto_cantidad();
\end{lstlisting}

\begin{nota}
\textbf{Por qué \texttt{AFTER UPDATE} y no \texttt{BEFORE}.} La
auditoría sólo tiene sentido cuando la modificación ya ocurrió. Si
algún \texttt{CHECK} o trigger anterior aborta el \texttt{UPDATE}, no
debe haber registro de auditoría: usando \texttt{AFTER UPDATE},
PostgreSQL y MySQL garantizan que el trigger sólo dispara cuando la
operación es exitosa.
\end{nota}

\subsection{c) Autonumerado de \texttt{parquimetro}}

Sobre la base \texttt{practico1b} (estacionamiento del Práctico 1.b)
hay que hacer que \texttt{id\_parquimetro} se asigne automáticamente.

\subsubsection*{MySQL 8+}

\begin{lstlisting}
USE practico1b_mysql;

ALTER TABLE parquimetro
    MODIFY id_parquimetro INT NOT NULL AUTO_INCREMENT;
\end{lstlisting}

Después de ejecutar esto, basta con omitir la columna en el
\texttt{INSERT}:

\begin{lstlisting}
INSERT INTO parquimetro (calle, altura) VALUES ('Colon', 1500);
SELECT * FROM parquimetro ORDER BY id_parquimetro;
\end{lstlisting}

\subsubsection*{Oracle XE 10/11 (\texttt{SEQUENCE + TRIGGER})}

Oracle 10/11 XE no tiene \texttt{IDENTITY} (apareció en Oracle 12c),
así que el patrón clásico es \emph{secuencia + trigger}. La
secuencia genera valores únicos crecientes; el trigger los asigna a
\texttt{:NEW.id\_parquimetro} cuando la sentencia no provee uno
explícito.

\begin{lstlisting}
-- Limpieza opcional para re-ejecutar el script.
BEGIN EXECUTE IMMEDIATE 'DROP TRIGGER bi_parquimetro_autonumerado';
EXCEPTION WHEN OTHERS THEN NULL; END;
/
BEGIN EXECUTE IMMEDIATE 'DROP SEQUENCE seq_parquimetro';
EXCEPTION WHEN OTHERS THEN NULL; END;
/

-- Crear la secuencia arrancando despues del MAX existente,
-- para no chocar con datos cargados.
DECLARE
    v_start NUMBER;
BEGIN
    SELECT NVL(MAX(id_parquimetro), 0) + 1
      INTO v_start
      FROM parquimetro;

    EXECUTE IMMEDIATE
      'CREATE SEQUENCE seq_parquimetro START WITH ' || v_start ||
      ' INCREMENT BY 1 NOCACHE NOCYCLE';
END;
/

CREATE OR REPLACE TRIGGER bi_parquimetro_autonumerado
BEFORE INSERT ON parquimetro
FOR EACH ROW
WHEN (NEW.id_parquimetro IS NULL)
BEGIN
    SELECT seq_parquimetro.NEXTVAL
      INTO :NEW.id_parquimetro
      FROM dual;
END;
/
\end{lstlisting}

Prueba:

\begin{lstlisting}
INSERT INTO parquimetro (calle, altura) VALUES ('Colon', 1500);
SELECT * FROM parquimetro ORDER BY id_parquimetro;
\end{lstlisting}

\paragraph{Comparación con MySQL.} En MySQL es suficiente declarar la
columna como \texttt{AUTO\_INCREMENT}. En Oracle 10/11 XE el mismo
efecto requiere dos objetos cooperando: la secuencia
(\texttt{seq\_parquimetro}) que produce valores y el trigger
(\texttt{bi\_parquimetro\_autonumerado}) que los inyecta en la fila
sólo cuando el cliente no provee uno explícito. La cláusula
\texttt{WHEN (NEW.id\_parquimetro IS NULL)} evita pisar los inserts
con ID explícito --- típicos en cargas de datos de prueba.

% =====================================================
\section*{Conclusión}
\addcontentsline{toc}{section}{Conclusión}
% =====================================================

Los cinco ejercicios cubren el espectro de programación del lado
servidor: \emph{triggers} para reglas que dependen de varias tablas
(Ej.\ 1), \emph{funciones} con cálculo manual y parámetros
\texttt{IN}/\texttt{OUT} (Ej.\ 2), \emph{procedimientos} con
validación, transaccionalidad y mensajes por pantalla (Ej.\ 3),
\emph{procedimientos} con seguridad por invocador (Ej.\ 4) y
\emph{triggers} de normalización, auditoría y autonumerado (Ej.\ 5).

Los puntos importantes para llevarse al parcial son:

\begin{itemize}
    \item \textbf{BEFORE/AFTER + INSERT/UPDATE/DELETE} es la matriz
    básica del trigger; \texttt{NEW} y \texttt{OLD} dan acceso a la
    fila nueva y vieja.
    \item Para \emph{abortar} desde el trigger: \texttt{SIGNAL
    SQLSTATE '45000'} en MySQL, \texttt{RAISE EXCEPTION} en
    PostgreSQL, \texttt{RAISE\_APPLICATION\_ERROR} en Oracle.
    \item Las \emph{funciones} se asocian a \texttt{RETURN} de un valor;
    los \emph{procedimientos} se invocan con \texttt{CALL} (MySQL,
    PostgreSQL) o en bloques \texttt{BEGIN \dots\ END} (Oracle) y
    pueden tener parámetros \texttt{IN}, \texttt{OUT} e \texttt{INOUT}.
    \item Para \emph{auditar el usuario de aplicación} se usan
    variables de sesión: \texttt{@var} en MySQL,
    \texttt{current\_setting(\dots)} en PostgreSQL.
    \item \texttt{SQL SECURITY INVOKER}/\texttt{DEFINER} en MySQL es la
    pieza que conecta el Práctico 3 con el Práctico 2 (DCL): el
    primero asume permisos del invocador, el segundo (default) los del
    creador.
    \item \texttt{AUTO\_INCREMENT} (MySQL) y \texttt{IDENTITY}
    (PostgreSQL/Oracle 12+) ahorran código; en Oracle 10/11 XE el
    patrón es \texttt{SEQUENCE + TRIGGER BEFORE INSERT}.
\end{itemize}

\end{document}
